% Generated draft: requires mathematical and source review.
\documentclass[11pt]{article}
\usepackage[margin=25mm]{geometry}
\usepackage{fontspec}
\setmainfont{Latin Modern Roman}
\usepackage{amsmath,amssymb,mathtools,unicode-math}
\setmathfont{Latin Modern Math}
\usepackage{enumitem,needspace,xurl,hyperref}
\setlength{\parindent}{0pt}
\setlength{\parskip}{5pt}
\setlength{\emergencystretch}{3em}
\newcommand{\N}{\mathbb{N}}
\newcommand{\Z}{\mathbb{Z}}
\newcommand{\Q}{\mathbb{Q}}
\newcommand{\R}{\mathbb{R}}
\newcommand{\C}{\mathbb{C}}
\newcommand{\F}{\mathbb{F}}
\newcommand{\A}{\mathbb{A}}
\newcommand{\PP}{\mathbb{P}}
\newcommand{\E}{\mathbb{E}}
\newcommand{\eps}{\varepsilon}
\newcommand{\dd}{\,\mathrm{d}}
\begin{document}
\section*{UnsolvedMath ID: 800657}
\textbf{Source problem code:} catalogue entry
\subsection{Definitions and mathematical statement}
For interpreting limits of sequences indexed by \(\N\), let \(M\) be a collection of ultrafilters on \(\N\), where an ultrafilter is a maximal proper filter. For a sequence \((a_n)_{n\in\N}\) in a topological space \(Y\), write
\[
p\text{-}\lim_{n\in\N}a_n=y
\quad\Longleftrightarrow\quad
\{n\in\N:a_n\in U\}\in p
\text{ for every neighborhood }U\text{ of }y.
\]
The supplied record does not specify whether the ultrafilters must be free or whether \(M\) must be nonempty.

Following the definitions attributed in the record to Kombarov, \(Y\) is \emph{weakly \(M\)-sequential} if
\[
\forall A\subseteq Y,\quad
A\ne\overline A
\ \Longrightarrow\
\exists y\in Y\setminus A\;
\exists(a_n)_{n\in\N}\in A^{\N}\;
\exists p\in M:\quad
p\text{-}\lim_{n\in\N}a_n=y.
\]
It is \emph{strongly \(M\)-sequential} if instead
\[
\forall A\subseteq Y,\quad
A\ne\overline A
\ \Longrightarrow\
\exists y\in Y\setminus A\;
\exists(a_n)_{n\in\N}\in A^{\N}\;
\forall p\in M:\quad
p\text{-}\lim_{n\in\N}a_n=y.
\]
Thus, in the strong version, the same point and sequence must work for every \(p\in M\).

Under the usual real-valued convention, \(C_p(X)\) denotes the set of continuous functions \(X\to\R\), with the subspace topology inherited from the product \(\R^X\). The supplied record does not explicitly specify this convention or the class of spaces \(X\).

\textbf{Question (source formulation).} Assuming that \(X\) satisfies the \(\delta_M\) property, must \(C_p(X)\) be weakly \(M\)-sequential? Equivalently, under the preceding convention for \(C_p(X)\), does the assumption imply that every nonclosed \(A\subseteq C_p(X)\) admits \(f\in C_p(X)\setminus A\), a sequence \((f_n)_{n\in\N}\) in \(A\), and \(p\in M\) such that
\[
\forall x\in X\;\forall\varepsilon>0,\qquad
\{n\in\N:|f_n(x)-f(x)|<\varepsilon\}\in p?
\]
\textbf{Unresolved qualification.} The supplied record does not define the \(\delta_M\) property or specify its admissible spaces and ultrafilter collections. Consequently, the hypothesis above remains an undefined source predicate, and the question cannot be made fully self-contained from the supplied evidence.
\subsection{Short English statement}
The question asks whether the source's \(\delta_M\) property of \(X\) forces every nonclosed subset of \(C_p(X)\) to contain a sequence with a limit outside that subset along some ultrafilter in \(M\). The definition of \(\delta_M\), the permitted spaces, and the restrictions on \(M\) are missing. The real-valued interpretation of \(C_p(X)\) requires confirmation.
\subsection{Sources}
\begin{itemize}
\item Local UnsolvedMath registry: permanent ID \(800657\), exact source code \texttt{catalogue entry}, supplied complete source record; upstream revision \texttt{d371ee917ac9647fbb06cd6fc529deec6f3dda82}. The registry attributes the problem to Tsaban, gives the year 2003, and labels it open. This is a supplied status label, not an independently verified mathematical or literature assessment.
\item Tsaban, \emph{Selection Principles in Mathematics: A Milestone of Open Problems} (2003), Problem 30. Supplied source URL: \url{https://arxiv.org/abs/math/0312182}. The record reports access on 2026-10-16 and extraction by \texttt{arxiv\_latex}; these are supplied provenance details, not a claim of independent source access.
\item Supplied attribution/context: [GTpseq, Problem 0.6 (reformulated)]. The full bibliographic reference and the original formulation are not supplied.
\item The weak and strong \(M\)-sequentiality definitions are attributed to Kombarov [Km]. The full bibliographic reference for [Km] is not supplied.
\end{itemize}
\end{document}
